\chapter{Ringkasan Kode Sumber}


Lampiran ini berfungsi sebagai penunjuk menuju repositori sumber yang digunakan dalam penelitian. Seluruh potongan kode RTL dan firmware yang relevan telah dijelaskan pada Bab IV, sehingga lampiran ini tidak mengulang listing kode secara lengkap.

\section{Repositori}

Seluruh kode sumber, dokumentasi proyek, berkas Vivado, firmware pengujian, serta skrip pendukung yang digunakan dalam penelitian ini tersedia secara terbuka pada repositori GitHub berikut:

\begin{center}
\url{https://github.com/AdakHaddad/thesi2s}
\end{center}

Repositori tersebut mencakup implementasi IP I2S Transmitter berbasis AXI4-Lite, firmware pengujian pada MicroBlaze V, berkas konfigurasi Vivado, skrip rekonstruksi proyek, serta dokumentasi tambahan yang digunakan selama proses pengembangan dan validasi sistem.

\section{Kode Testbench}

Testbench berikut digunakan untuk memverifikasi transaksi AXI4-Lite, konfigurasi
register IP I2S, serta keluaran clock dan sinyal audio I2S. Testbench bersifat
\textit{self-checking} dan mencatat jumlah kesalahan selama simulasi.

\lstinputlisting[
    language=Verilog,
    caption={Testbench AXI4-Lite dan protokol I2S.},
    label={lst:i2s-testbench},
    breaklines=true,
    basicstyle=\ttfamily\scriptsize
]{../freshi2s/project_1.srcs/sim_1/new/new_tb.v}

